% NISQ-Folio paper draft. Build: pdflatex main && bibtex main && pdflatex main && pdflatex main
% (or upload the paper/ folder to Overleaf). Tables in generated/ come from paper/make_tables.py.
% Text marked \pending{...} depends on the Week 12 noisy grid and must be completed before submission.
\documentclass[11pt]{article}
\usepackage[a4paper,margin=2.5cm]{geometry}
\usepackage[T1]{fontenc}
\usepackage[utf8]{inputenc}
\usepackage{lmodern}
\usepackage{amsmath,amssymb,amsthm}
\usepackage{booktabs}
\usepackage{graphicx}
\usepackage[numbers,sort&compress]{natbib}
\usepackage{xcolor}
\usepackage[colorlinks=true,linkcolor=blue!50!black,citecolor=blue!50!black,urlcolor=blue!50!black]{hyperref}

\newcommand{\pending}[1]{\textcolor{red}{[PENDING (Week 12 noisy grid): #1]}}
\newcommand{\todo}[1]{\textcolor{orange}{[TODO: #1]}}
\newcommand{\Amp}{\mathcal{A}}

\title{Constraint-preserving and adaptive mixers in noisy QAOA for cardinality-constrained
portfolio selection: a pre-registered benchmark against classical baselines}

\author{
  Yuvaraj Guliani\thanks{Equal contribution.}\\
  S P Jain School of Global Management\\
  \texttt{yuvaraj.bs24dmu057@spjain.org}
  \and
  Kallol Das\footnotemark[1]\\
  \todo{affiliation and e-mail}
}
\date{Draft, \today}

\begin{document}
\maketitle

\begin{abstract}
We benchmark a hybrid quantum--classical pipeline for cardinality-constrained portfolio
selection: random-matrix-theory (RMT) covariance denoising, a QUBO encoding with a provably
exact cardinality penalty, and the quantum approximate optimization algorithm (QAOA) with a
standard transverse-field mixer, a constraint-preserving XY mixer, and an adaptive (ADAPT)
mixer. On 30 US equities (2015--2023) and problem sizes of 8 to 14 assets, every design choice,
tuning rule and analysis was pre-registered before the corresponding results existed, and all
comparisons are paired over the same 30 out-of-sample instances per cell. In noiseless
simulation, the XY and ADAPT mixers place 2.8--3.1$\times$ and 4.0--4.9$\times$ more probability
on the optimal portfolio than uniform random feasible guessing, whereas the penalized standard
mixer does not exceed random guessing at any size. Classical simulated annealing nevertheless
finds the optimum on at least 90\% of instances with 50--200 objective evaluations, two to four
orders of magnitude fewer samples than QAOA uses. RMT denoising improves neither the classical
optimum out of sample nor QAOA. Under a depolarizing-plus-readout noise model calibrated to IBM
Heron-class median error rates, \pending{the noise level at which each mixer's advantage over
random feasible guessing vanishes, with confidence intervals, and how it scales with $N$}.
Code, configurations, run records and the full decision log are openly available.
\end{abstract}

\section{Introduction}
Portfolio selection under a cardinality constraint (choosing exactly $K$ of $N$ assets) is a
natural target for near-term quantum optimization: the mean-variance objective of
\citet{markowitz1952} becomes a quadratic unconstrained binary optimization (QUBO) problem once
the constraint is moved into a penalty~\citep{lucas2014}, and QUBOs map directly onto the cost
Hamiltonians of the quantum approximate optimization algorithm (QAOA)~\citep{farhi2014qaoa}.
Portfolio optimization is among the most studied financial applications of quantum
computing~\citep{egger2020,herman2023}, including QAOA studies with hard constraints and
finite-shot and hardware noise~\citep{hodson2019,brandhofer2022}.

Two questions are central for noisy intermediate-scale quantum (NISQ) devices~\citep{preskill2018}.
First, does a carefully engineered pipeline beat \emph{classical} baselines? Second, how much
noise can it tolerate before any quantum effect disappears? General bounds suggest that noise
strongly limits optimization on current devices~\citep{stilckfranca2021}; we ask the question
empirically for a specific, complete pipeline.

We combine three ingredients that are usually studied separately: (i) RMT eigenvalue clipping of
the covariance matrix~\citep{laloux1999,bun2017}; (ii) a QUBO with a cardinality penalty that is
\emph{provably} exact without knowing the optimum; and (iii) QAOA with the standard X mixer, the
constraint-preserving XY mixer of the quantum alternating operator ansatz~\citep{hadfield2019,wang2020xy}
started from a Dicke state~\citep{bartschi2019dicke}, and the adaptive ADAPT-QAOA mixer~\citep{zhu2022adapt}.

\paragraph{Contributions.}
\begin{itemize}
  \item A fully pre-registered, paired factorial benchmark (denoising $\times$ mixer $\times$ noise
        level $\times$ problem size, 30 instances per cell), with every decision recorded before the
        corresponding result existed~\citep{nosek2018}.
  \item A success metric that cannot be met by blind sampling: amplification over random feasible
        guessing, measured exactly on the final output distribution (Section~\ref{sec:metric}).
  \item Evidence that constraint-preserving and adaptive mixers, but not the penalized standard
        mixer, beat random feasible guessing in noiseless simulation, while classical heuristics
        remain orders of magnitude cheaper at these sizes.
  \item A null result for RMT denoising on this universe, with a mechanism.
  \item \pending{the measured noise crossover per mixer and its dependence on $N$}.
\end{itemize}

\section{Problem formulation}\label{sec:problem}
Let $\mu\in\mathbb{R}^N$ and $\Sigma\in\mathbb{R}^{N\times N}$ be the annualized mean returns and
covariance of $N$ assets, estimated from daily returns. A selection $x\in\{0,1\}^N$ with
$\mathbf{1}^\top x=K$ is held in equal weights $w=x/K$, and we minimize the mean-variance objective
\begin{equation}
  f(x) \;=\; \lambda\, w^\top\Sigma w - \mu^\top w
       \;=\; \frac{\lambda}{K^2}\, x^\top\Sigma x - \frac{1}{K}\,\mu^\top x ,
  \qquad \mathbf{1}^\top x = K,
  \label{eq:objective}
\end{equation}
with $K=N/2$. Equal weights make the decision purely binary, so one qubit encodes one asset.

\paragraph{Risk aversion.} For each instance we compute
$\lambda_{\mathrm{bal}}=\mathrm{sd}_S[\mu^\top w]/\mathrm{sd}_S[w^\top\Sigma w]$ over all
$\binom{N}{K}$ selections $S$, the value at which return and risk vary equally across the feasible
set. We fix $\lambda=4.5$, the median of $\lambda_{\mathrm{bal}}$ over 200 seeded in-sample
instances rounded to $0.5$, chosen before any solver was evaluated. Smaller $\lambda$ makes the
problem collapse towards ``pick the $K$ highest-return assets''.

\paragraph{QUBO and an exact penalty.} With $a=\lambda/K^2$, $b=1/K$ and penalty weight $P$,
\begin{equation}
  E(x) = f(x) + P(\mathbf{1}^\top x - K)^2
       = x^\top\!\left(a\Sigma + P\mathbf{1}\mathbf{1}^\top + \operatorname{diag}(-b\mu - 2PK)\right)\! x + PK^2 .
\end{equation}
Let $f_{\mathrm{ub}}$ be the objective of any feasible selection (we use greedy) and, for each
cardinality $m$, $L_m = a\,\lambda_{\min}(\Sigma)\,m - b\sum_{i\le m}\mu_{(i)}$, where $\mu_{(i)}$
are the returns in decreasing order. Then $L_m\le f(x)$ for every $x$ with $m$ ones, and choosing
\begin{equation}
  P > \max_{m\ne K}\; \frac{f_{\mathrm{ub}} - L_m}{(m-K)^2}
  \label{eq:penalty}
\end{equation}
guarantees $E(x) > f_{\mathrm{ub}} \ge f^\ast$ for every infeasible $x$, so the global QUBO minimum
is the constrained optimum. The bound uses only cheap quantities; it never uses $f^\ast$. On real
instances it is exact in 200 of 200 cases and 7--11 times larger than the smallest exact penalty.

\section{Methods}\label{sec:methods}
\subsection{Data and covariance estimation}
We use daily adjusted closes of 30 large US equities, five in each of six GICS sectors, from
2015-01-02 to 2023-12-29 (2{,}264 trading days, no missing values). Moments are estimated on
2015--2020 (1{,}510 days); selected portfolios are evaluated on 2021--2023, which no solver sees.
The selection is subject to survivorship bias: all names were chosen knowing they existed through
2023, so absolute returns are optimistic; all solvers see identical inputs.

\paragraph{RMT denoising.} For $N$ standardized uncorrelated series of length $T$, sample
correlation eigenvalues follow the Marchenko--Pastur law~\citep{marchenko1967} on
$[\lambda_-,\lambda_+]=[(1-\sqrt{q})^2,(1+\sqrt{q})^2]$, $q=N/T$. Following~\citet{laloux1999}
and~\citet{bun2017}, eigenvalue clipping keeps eigenvalues above $\lambda_+$, replaces the rest by
their mean (preserving the trace), keeps the eigenvectors, and renormalizes to unit diagonal. We
clean the 30-asset universe correlation ($q=0.020$) and each instance uses its block. We also
report Ledoit--Wolf shrinkage~\citep{ledoit2004} as a non-RMT reference.

\subsection{QAOA and mixers}
For cost Hamiltonian $H_C$ (the Ising form of $E$ or $f$, rescaled by the standard deviation of its
energy under uniform spins) and mixer $M$, the depth-$p$ state is
$\lvert\gamma,\beta\rangle=\prod_{l=1}^{p} e^{-i\beta_l M_l} e^{-i\gamma_l H_C}\lvert\psi_0\rangle$.
\begin{itemize}
  \item \textbf{X mixer:} $M=\sum_i X_i$, $\lvert\psi_0\rangle=\lvert{+}\rangle^{\otimes N}$, cost $E$ with
        penalty $P=c\times$\eqref{eq:penalty}.
  \item \textbf{XY ring:} $M$ is one parity-ordered Trotter step of $\sum_i (X_iX_{i+1}+Y_iY_{i+1})/2$,
        which preserves Hamming weight~\citep{hadfield2019,wang2020xy}; $\lvert\psi_0\rangle$ is the Dicke
        state $\lvert D^N_K\rangle$~\citep{bartschi2019dicke}; cost $f$ only, since infeasible states are
        never reached noiselessly.
  \item \textbf{ADAPT:} following~\citet{zhu2022adapt}, layer $l$'s mixer is chosen from a pool of
        weight-preserving operators $\{$XY ring$\}\cup\{(X_iX_j+Y_iY_j)/2\}$ by the largest shot-estimated
        energy gradient; the shots spent on this selection are counted in every budget comparison.
\end{itemize}
Each energy evaluation uses $S=1{,}024$ shots sampled from the exact distribution, as on hardware;
angles are optimized with COBYLA~\citep{powell1994} or SPSA~\citep{spall1992} within 100
evaluations, from a linear-ramp initialization. Settings (depth, optimizer, penalty scale) were
chosen on separate tuning instances by a pre-registered rule; the selections are X: COBYLA,
$p=2$, $c=0.1$; XY ring: SPSA, $p=1$; ADAPT: COBYLA, $p=4$.

\subsection{Success metric}\label{sec:metric}
A QAOA run draws $10^4$--$10^6$ samples, more than the $2^N\le16{,}384$ bitstrings at our sizes,
so the best sampled solution would also be found by blind guessing. We therefore score the
final distribution $p(x)$ exactly by its amplification over uniform random feasible guessing,
\begin{equation}
  \Amp = \frac{p_{\mathrm{opt}}}{m/\binom{N}{K}}, \qquad p_{\mathrm{opt}} = \sum_{x\ \text{optimal}} p(x),
\end{equation}
where $m$ is the number of optimal selections. $\Amp>1$ means QAOA beats random feasible guessing,
and probability on infeasible states counts against it. The \emph{noise crossover} is the noise
level at which the median $\Amp$ across instances falls to $1$. We also report the expected
normalized gap, whether the modal (most probable) feasible selection is optimal, and the
out-of-sample risk of that modal selection.

\subsection{Noise model and simulation}
A single knob, the noise level $m_{\mathrm{noise}}$, multiplies the median error rates of the IBM
Heron-class fake backend FakeKingston (single-qubit $2.36\times10^{-4}$, CZ $1.83\times10^{-3}$, readout
$9.52\times10^{-3}$), so $m_{\mathrm{noise}}=1$ approximates today's hardware. Gate errors are
depolarizing with parameter $r\,d/(d-1)$ ($d=2^n$), so that the average gate infidelity equals the
error rate $r$; readout is an independent symmetric bit flip. Circuits (including a Dicke-state
preparation and XY gates) are transpiled onto a connected $N$-qubit region of the device's
heavy-hex coupling map and simulated exactly as density matrices with Qiskit and Qiskit
Aer~\citep{javadiabhari2024qiskit}. At $m_{\mathrm{noise}}=0$ the noisy pipeline reproduces the
noiseless simulator to $10^{-15}$. Thermal relaxation and crosstalk are not modelled.

\subsection{Classical baselines}
Exhaustive search (the exact optimum), greedy forward selection, simulated annealing~\citep{kirkpatrick1983}
with feasibility-preserving swap moves at several step budgets, and the continuous long-only
Markowitz relaxation rounded to its $K$ largest weights.

\section{Experimental design}\label{sec:design}
A cell is (denoising $\in\{$none, RMT$\}$, mixer, noise level, $N$). Every cell uses the same 30
evaluation instances per $N$ (random $N$-subsets of the universe, drawn from a seeded stream),
so all comparisons are paired. Noise levels $m_{\mathrm{noise}}\in\{0,0.25,0.5,1\}$ were fixed by a
pre-registered pilot on separate tuning instances; noisy runs cover $N\in\{8,10,12\}$ (ADAPT:
$N\le10$, as exact noisy ADAPT at $N=12$ is computationally infeasible here), and $N=14$ is
noiseless only. The full grid has 2{,}400 tasks. Every run records its configuration, data
checksum and code commit; all randomness derives from one root seed. Paired comparisons use
two-sided Wilcoxon signed-rank tests on $\ln\Amp$ differences; because instances share assets,
$p$-values are descriptive.

\section{Results}\label{sec:results}
\subsection{Classical baselines}
Table~\ref{tab:classical} shows that classical heuristics are strong at these sizes: greedy is
optimal on 87--93\% of instances and simulated annealing reaches 100\% with 500 steps.
\begin{table}[ht]\centering\small
\caption{Share of evaluation instances solved optimally (sample covariance, 30 instances per $N$).}
\label{tab:classical}
% Generated by paper/make_tables.py from results/experiments/main_r2; commits 31c7e80; 2026-10-20T21:52:38+00:00
\begin{tabular}{lcccc}
\toprule
Method & $N=8$ & $N=10$ & $N=12$ & $N=14$ \\
\midrule
Greedy & 87\% & 93\% & 90\% & 87\% \\
Markowitz + rounding & 83\% & 77\% & 57\% & 33\% \\
SA, 25 steps & 67\% & 40\% & 7\% & 20\% \\
SA, 50 steps & 90\% & 90\% & 57\% & 27\% \\
SA, 100 steps & 100\% & 97\% & 67\% & 60\% \\
SA, 200 steps & 100\% & 100\% & 100\% & 90\% \\
SA, 500 steps & 100\% & 100\% & 100\% & 100\% \\
SA, 2000 steps & 100\% & 100\% & 100\% & 100\% \\
\bottomrule
\end{tabular}

\end{table}

\subsection{Noiseless QAOA}
\begin{table}[ht]\centering\small
\caption{Noiseless amplification $\Amp$ over random feasible guessing: median [IQR] (share of
instances with $\Amp>1$), sample covariance.}
\label{tab:amp}
\resizebox{\linewidth}{!}{% Generated by paper/make_tables.py from results/experiments/main_r2; commits 31c7e80; 2026-10-20T21:52:38+00:00
\begin{tabular}{lcccc}
\toprule
Mixer & $N=8$ & $N=10$ & $N=12$ & $N=14$ \\
\midrule
ADAPT & 4.01 [3.07, 6.13] (93\%) & 4.91 [3.84, 8.33] (100\%) & 4.92 [3.13, 8.58] (100\%) & 4.39 [3.36, 7.95] (100\%) \\
XY ring & 2.87 [1.94, 3.80] (93\%) & 3.08 [2.45, 5.39] (90\%) & 2.82 [2.03, 4.02] (90\%) & 3.08 [1.26, 3.99] (77\%) \\
X & 0.73 [0.65, 0.85] (20\%) & 0.81 [0.70, 0.93] (17\%) & 0.90 [0.73, 1.06] (37\%) & 0.77 [0.66, 0.88] (13\%) \\
\bottomrule
\end{tabular}
}
\end{table}
By the pre-registered criterion (median $\Amp>1$ and $\Amp>1$ on at least two thirds of
instances), the XY and ADAPT mixers beat random feasible guessing at every size and the X mixer
at none (Table~\ref{tab:amp}). Paired over instances, ADAPT exceeds XY by a geometric-mean factor
of 1.6--2.3 and XY exceeds X by 2.9--3.5 (Table~\ref{tab:mixers}). The X mixer's weakness tracks
the penalty: smaller penalty scales were consistently better in tuning, and a shot-noise-robust
optimizer (SPSA) did not help it.
\begin{table}[ht]\centering\small
\caption{Paired mixer comparisons: geometric-mean ratio of $\Amp$ (share of instances where the
first mixer is better).}
\label{tab:mixers}
% Generated by paper/make_tables.py from results/experiments/main_r2; commits 31c7e80, bb8eca2; 2026-10-22T02:44:56+00:00
\begin{tabular}{lcccc}
\toprule
Comparison & $N=8$ & $N=10$ & $N=12$ & $N=14$ \\
\midrule
ADAPT / XY ring & 1.61 (80\%) & 1.61 (83\%) & 1.65 (97\%) & 2.29 (90\%) \\
XY ring / X & 3.47 (97\%) & 3.44 (93\%) & 2.92 (93\%) & 2.86 (87\%) \\
ADAPT / X & 6.53 (93\%) & 6.08 (100\%) & 5.11 (100\%) & 7.45 (100\%) \\
\bottomrule
\end{tabular}

\end{table}

\paragraph{Probability on the optimum versus the most likely answer.} A high $\Amp$ does not imply
that the most probable selection is optimal (Table~\ref{tab:deployment}): the X mixer's modal
feasible selection is optimal on 63--77\% of instances despite $\Amp<1$, while the
constraint-preserving mixers spread probability over many good selections; their modal selections
are rarely optimal, but their realized out-of-sample variance is only a few percent above the
optimum's.
\begin{table}[ht]\centering\small
\caption{Modal feasible selection: share optimal / median ratio of its realized 2021--2023 variance
to the optimum's (sample covariance).}
\label{tab:deployment}
% Generated by paper/make_tables.py from results/experiments/main_r2; commits 31c7e80; 2026-10-20T21:52:38+00:00
\begin{tabular}{lcccc}
\toprule
Mixer & $N=8$ & $N=10$ & $N=12$ & $N=14$ \\
\midrule
ADAPT & 23\% / 1.05 & 7\% / 1.02 & 0\% / 1.05 & 3\% / 1.12 \\
XY ring & 10\% / 1.03 & 3\% / 1.05 & 3\% / 1.05 & 0\% / 1.16 \\
X & 77\% / 1.00 & 77\% / 1.00 & 77\% / 1.00 & 63\% / 1.00 \\
\bottomrule
\end{tabular}

\end{table}

\subsection{Sampling budget}
Table~\ref{tab:budget} compares QAOA's sampling cost with the number of simulated-annealing steps
needed for 90\% optimal. Although a shot and an annealing step are different units, classical
search needs two to four orders of magnitude fewer samples at $N\le14$.
\begin{table}[ht]\centering\small
\caption{Median QAOA shots per run (optimization plus ADAPT operator selection) versus the smallest
simulated-annealing budget that is optimal on $\geq90\%$ of instances.}
\label{tab:budget}
% Generated by paper/make_tables.py from results/experiments/main_r2; commits 31c7e80, bb8eca2; 2026-10-22T02:44:56+00:00
\begin{tabular}{lcccc}
\toprule
Method & $N=8$ & $N=10$ & $N=12$ & $N=14$ \\
\midrule
ADAPT (shots) & 339,968 & 479,232 & 651,264 & 856,064 \\
XY ring (shots) & 102,400 & 102,400 & 102,400 & 102,400 \\
X (shots) & 40,960 & 40,960 & 37,888 & 38,400 \\
SA (steps, $\geq 90\%$ optimal) & 50 & 50 & 200 & 200 \\
\bottomrule
\end{tabular}

\end{table}

\subsection{RMT denoising}
On the universe, 25 of 30 correlation eigenvalues lie \emph{below} the Marchenko--Pastur band, which
noise cannot produce at $q=0.02$; clipping replaces 27 of 30 eigenvalues and effectively imposes a
three-factor model. When it changes the classical optimum (26--44\% of instances), the selected
portfolio has higher realized out-of-sample variance more often than not
(Table~\ref{tab:denoise}), and its correlation estimate is further from the out-of-sample
correlation than the sample estimate. In a synthetic one-factor model, clipping beat the sample
estimate on 10 of 10 seeds when the non-market spectrum was flat, but lost on 9 of 10 when genuine
structure lay below the band, the regime the real data resembles. Denoising also has no consistent
effect on QAOA (Table~\ref{tab:denoise-qaoa}).
\begin{table}[ht]\centering\small
\caption{RMT clipping versus sample covariance for the exact classical optimum, evaluated out of
sample (50 instances per $N$).}
\label{tab:denoise}
% Generated by paper/make_tables.py from results/denoise_evaluation/20261016T045617Z_01159b2.json; commits 01159b2; 2026-10-22T02:44:56+00:00
% correlation distance to out-of-sample: sample 5.643, rmt_clip 5.943, ledoit_wolf 5.329
\begin{tabular}{cccc}
\toprule
$N$ & Selection changed & Lower realized variance (among changed) & Wilcoxon $p$ \\
\midrule
8 & 13/50 & 46\% & 0.455 \\
10 & 15/50 & 53\% & 0.389 \\
12 & 22/50 & 27\% & 0.00744 \\
14 & 18/50 & 44\% & 0.196 \\
\bottomrule
\end{tabular}

\end{table}
\begin{table}[ht]\centering\small
\caption{Effect of RMT denoising on QAOA: paired geometric-mean ratio of $\Amp$ (share of instances
improved).}
\label{tab:denoise-qaoa}
% Generated by paper/make_tables.py from results/experiments/main_r2; commits 31c7e80, bb8eca2; 2026-10-22T02:44:56+00:00
\begin{tabular}{lcccc}
\toprule
Comparison & $N=8$ & $N=10$ & $N=12$ & $N=14$ \\
\midrule
ADAPT: RMT / sample & 1.09 (60\%) & 1.04 (53\%) & 0.79 (40\%) & 1.12 (53\%) \\
XY ring: RMT / sample & 1.02 (53\%) & 1.01 (53\%) & 0.95 (47\%) & 1.11 (60\%) \\
X: RMT / sample & 0.89 (17\%) & 0.98 (43\%) & 0.87 (30\%) & 0.93 (47\%) \\
\bottomrule
\end{tabular}

\end{table}

\subsection{Noise}
\pending{amplification versus noise level per mixer and $N$ (figure); crossover estimates with
bootstrap confidence intervals; effect on feasibility leakage and out-of-sample risk; comparison of
the $N=8$ results under full-device and compact routing as a robustness check}.

\section{Discussion and limitations}
At the sizes that can be simulated exactly, the quantum part of the pipeline shows a genuine but
modest effect, and only with mixers that respect the constraint. It does not compete with classical
heuristics on cost, and its first stage (RMT denoising) does not help on this universe.
\pending{interpretation of the noise crossover relative to current hardware}.

Limitations: equal weights rather than continuous weights; a single 30-stock universe with
survivorship bias and a single out-of-sample period; $N\le14$, where exact search is trivial, so
conclusions concern trends rather than practical speed-ups; a depolarizing-plus-readout noise model
without relaxation or crosstalk, applied uniformly to all qubits; exact noisy ADAPT only up to
$N=10$; and descriptive rather than confirmatory $p$-values.

\section{Conclusion}
\pending{one-paragraph conclusion once the noise crossover is measured}.

\section*{Data and code availability}
All code, configurations, the decision log (27 numbered decision records, including pre-registrations and
their results), and the scripts that generate every table are available at
\url{https://github.com/yuvi1469/nisq-folio}. Raw market data are not redistributed; the code
downloads them from Yahoo Finance and verifies the snapshot by checksum.

\section*{Author contributions}
Both authors contributed equally to this work.

\section*{Use of AI assistance}
Software development, analysis scripts and drafting of this manuscript were assisted by an AI
system (Anthropic's Claude). All references were checked against their DOI or arXiv records. The
authors take full responsibility for the content. \todo{confirm this wording meets the target
journal's policy}

\bibliographystyle{unsrtnat}
\bibliography{references}
\end{document}
